To prepare the analytic argument, we now make the transform available for compactness and for tests that
need not be $L^2$ inputs on the whole support. We then fix the geometric
quantities whose packing will imply Theorem~\ref{thm:main}.
The ball-map conclusion has no instances when $D_E=0$.
We therefore work with positive-diameter support in the geometric
argument and restore the degenerate case in the final proof.
Constants denoted by $C$ may change between occurrences; their
dependencies are stated in each lemma or at the start of the argument.
We normalize $\H^n$ so that its restriction to an $n$-plane is
$n$-dimensional Lebesgue measure.

\subsection{Global growth and hard truncations}

\begin{lemma}\label{lem:global-growth}
Let $1\le n\le d$, and let $\mu$ be $n$-Ahlfors--David regular with
$D_E>0$. Then
\begin{equation}\label{eq:global-growth}
 \mu(B(a,r))\le G r^n\qquad(a\in\R^d,\ r>0),
 \qquad G=9^d C_{\rm AD}.
\end{equation}
If $D_E<\infty$, then
$\mu(\R^d)\le9^d C_{\rm AD}(D_E/2)^n$.
For every $\varepsilon>0$, every $f\in L^2(\mu)$, and every $a\in\R^d$,
the integral defining $R_{\mu,\varepsilon}f(a)$ is absolutely convergent.
\end{lemma}

\begin{proof}
A Radon measure on Euclidean space is concentrated on its support.
If $D_E<\infty$, choose a maximal $D_E/4$-separated subset of $E$.
The disjoint ambient balls of radius $D_E/8$ around its points lie in
a ball of radius $9D_E/8$ around any fixed point of $E$.
Comparison of their $d$-dimensional volumes bounds their number by
$9^d$. Maximality gives a cover of $E$ by closed balls of radius
$D_E/4$, hence by open balls of radius $D_E/2$. The asserted total-mass
bound follows from AD upper regularity at their centers.

If $\mu(B(a,r))>0$, choose $x\in E\cap B(a,r)$. When $2r\le D_E$,
$B(a,r)\subset B(x,2r)$ and the AD upper bound gives
$\mu(B(a,r))\le 2^n C_{\rm AD}r^n$. When $2r>D_E$, the preceding
total-mass estimate gives $\mu(B(a,r))\le9^dC_{\rm AD}r^n$.
For $D_E=\infty$ only the first case is needed. Since $n\le d$,
$2^n\le9^d$.

Finally, dyadic annuli and \eqref{eq:global-growth} give
\[
 \int_{|a-y|>\varepsilon}|K_n(a-y)|^2\,d\mu(y)
 \le 2^nG\varepsilon^{-n}\sum_{j=0}^{\infty}2^{-jn}<\infty.
\]
Cauchy--Schwarz proves the last assertion.
\end{proof}

In intermediate lemmas we use measures $\sigma$ with an explicit
global upper bound $G$ and a lower bound
\begin{equation}\label{eq:lower-growth}
 \sigma(B(x,r))\ge c r^n
 \qquad(x\in\supp\sigma,\quad 0<r\le D_{\supp\sigma}),
\end{equation}
where $c>0$. A \emph{fixed analytic class} means that
$d,n,G,c$ and a hard-truncation operator bound $D_0$ are fixed.
Sometimes the lower bound is assumed only inside a specified ball;
that local restriction will be stated explicitly. Measures with
\eqref{eq:global-growth} have no atoms when $n\ge1$.

\subsection{Mean-zero Riesz pairings}

For a real compactly supported Lipschitz function $\varphi$ satisfying
$\int\varphi\,d\sigma=0$, choose $a\in\R^d$ and $R>0$ such that
$\supp\varphi\subset B(a,R)$, and put $A=B(a,2R)$. Define the vector
pairing
\begin{align}
 \mathcal R_\sigma(\varphi)
 ={}&\frac12\iint_{A\times A}
 K_n(x-y)\bigl(\varphi(x)-\varphi(y)\bigr)
 \,d\sigma(x)d\sigma(y)\notag\\
 &+\int_A\int_{A^c}\varphi(x)
 \bigl(K_n(x-y)-K_n(a-y)\bigr)\,d\sigma(y)d\sigma(x).
 \label{eq:pairing-definition}
\end{align}
The value of the integrand on the diagonal is set to zero. A test in
a unit vector direction $e$ means $e\cdot\mathcal R_\sigma(\varphi)$.
For a vector test the pairing is the sum of its coordinate pairings,
with zero mean imposed on each coordinate. We call a measure
\emph{reflectionless} if this pairing vanishes for every compactly
supported Lipschitz test of zero measure mean.

\begin{lemma}\label{lem:pairings}
Suppose $\sigma$ satisfies global upper $n$-growth, with $n\ge1$.
The integrals in \eqref{eq:pairing-definition} converge absolutely,
and their sum is independent of $a,R$. For a fixed compact Lipschitz
test, the contribution of $|y-a|>T$ is $O(T^{-1})$ as $T\to\infty$.
If the hard truncations have $L^2$ bound $D_0$, smooth truncations with
a fixed bounded profile have bound at most $D_0+CG$. Pairings against
an indicator of finite measure have a limit as the inner truncation
tends to zero whenever that indicator equals one on a neighborhood of
the compact Lipschitz test's support.
\end{lemma}

\begin{proof}
The kernel is odd and satisfies
\[
 |K_n(z)|=|z|^{-n},\qquad |\nabla K_n(z)|\le C_n|z|^{-n-1}.
\]
For $h>0$, upper growth and a dyadic decomposition toward the origin
give
\begin{equation}\label{eq:near-growth-integral}
 \int_{0<|x-y|<h}|x-y|^{1-n}\,d\sigma(y)\le C_n G h.
\end{equation}
Multiplication by $\Lip\varphi$ bounds the diagonal singularity in
the first integral of \eqref{eq:pairing-definition}. The second
integral is separated from the support of $\varphi$ at its finite
boundary, and for $|y-a|\ge2R$ the kernel difference is bounded by
$C_nR|y-a|^{-n-1}$. Summation over exterior annuli proves absolute
convergence and the stated tail bound.

After enlarging to a common localization ball containing both
subtraction points, changing that point adds a vector independent of $x$,
whose integral against $\varphi$ is zero. Changing the localization
ball merely moves a finite cross term between the two integrals:
expanding it, using oddness, and again using $\int\varphi\,d\sigma=0$
shows that their sum is unchanged. This calculation can first be
made with both an inner and an outer cutoff; the estimates just proved
justify removing them.

The difference between a hard cutoff at $s$ and a bounded smooth
cutoff that vanishes on $[0,s]$ and equals one on $[2s,\infty)$ has
kernel bounded by $C s^{-n}\mathbf1_{\{s<|x-y|<2s\}}$.
Both its row and column integrals are bounded by $CG$. The Schur test
therefore gives the asserted $L^2$ bound. For an indicator input that
equals one on a neighborhood of the test support, antisymmetrization
on a smaller neighborhood has
the integrable majorant already established. Dominated convergence
identifies its limit with the local term in
\eqref{eq:pairing-definition}; separated exterior terms have ordinary
absolute convergence.
\end{proof}

For later use, if $\supp\varphi\subset B(a,r)$,
$\Lip\varphi\le r^{-1}$, and $T>4r$, the same proof gives
\begin{equation}\label{eq:localized-pairing-tail}
 \left|\mathcal R_\sigma(\varphi)
       -\langle R_\sigma\mathbf1_{B(a,T)},\varphi\rangle_\sigma\right|
 \le C G\frac rT\int|\varphi|\,d\sigma.
\end{equation}
Here the bracket denotes the limit of the truncated pairing, whose
existence follows from Lemma~\ref{lem:pairings}. It is only this
integrated limit that is used. The notation does not assert a
pointwise principal value.

\subsection{Weak limits and supports}

Local weak convergence of Radon measures means convergence of their
integrals against every compactly supported continuous function.
The following facts also apply to the translated and dilated measure
$\sigma_{a,r}(F)=r^{-n}\sigma(a+rF)$, for $a\in\R^d$ and $r>0$.

\begin{lemma}\label{lem:weak-compactness}
Let $\sigma_j$ be Radon measures with a common global upper growth
bound $G$. Suppose $0\in\supp\sigma_j$ and that a common lower bound
$c r^n$ holds at support centers up to radii $L_j\to\infty$.
There is a subsequence converging locally weakly to a nonzero measure
$\sigma$ with global upper and lower $n$-growth bounds.
The supports converge locally in the following sense:
\begin{enumerate}[label=(\roman*)]
\item each point of $\supp\sigma$ is a limit of points of
$\supp\sigma_j$;
\item any convergent sequence $x_j\in\supp\sigma_j$ has its limit
in $\supp\sigma$.
\end{enumerate}
Uniform hard-truncation $L^2$ bounds pass to $\sigma$. Moreover,
if $\varphi$ is compact Lipschitz with $\int\varphi\,d\sigma=0$,
it can be corrected to compact Lipschitz tests $\varphi_j$ with
$\int\varphi_j\,d\sigma_j=0$ so that
\[
 \mathcal R_{\sigma_j}(\varphi_j)\longrightarrow
 \mathcal R_\sigma(\varphi).
\]
The supports and Lipschitz constants of the corrected tests remain
uniformly bounded, and $\varphi_j\to\varphi$ uniformly.
\end{lemma}

\begin{proof}
Uniform mass bounds on compact balls give sequential compactness by
diagonal weak compactness on an exhaustion of $\R^d$. The upper
growth bound passes to the limit by testing slightly larger balls.
The lower bound at zero ensures nonzeroness. If $x\in\supp\sigma$,
every fixed neighborhood of $x$ has positive limiting mass and hence
meets the approximating supports eventually. If $x_j\to x$ and
$x_j\in\supp\sigma_j$, a fixed sufficiently small support ball at
$x_j$ has mass at least a fixed positive multiple of its radius to
the power $n$. A compactly supported bump near $x$, or the
Portmanteau inequality on a closed ball, transfers this positive mass
to the limit. This proves both support assertions and, by shrinking
and enlarging radii before taking limits, the lower growth bound.
Only fixed radii are used at this step, so they are less than $L_j$
eventually. Harmless dimensional changes in $c,G$ would also suffice.

We give the operator-bound argument to account for hard boundaries.
On a fixed pair of compact supports, the product measures
$\sigma_j\times\sigma_j$ converge weakly. For all but countably many
$\varepsilon>0$, the set $|x-y|=\varepsilon$ has zero product measure
under $\sigma\times\sigma$. At such a radius the truncated bilinear
form on a compact continuous scalar input and a compact continuous
$\R^d$-valued dual test passes to the limit. Their
$L^2$ norms pass to the limit as well, so its bound is the original
operator bound. Exhausting space leaves only a countable exceptional
set of radii. For any specified $\varepsilon>0$, choose good radii
decreasing to $\varepsilon$. Dominated convergence on compact supports,
away from the diagonal, gives the same bilinear estimate for the hard
cutoff $|x-y|>\varepsilon$. Density in $L^2$, together with the
$L^2$ integrability of the exterior kernel from
Lemma~\ref{lem:global-growth}, identifies the resulting bounded
operator with the actual integral for every $L^2$ input.

Choose a compact Lipschitz bump $\eta$ with $\int\eta\,d\sigma>0$.
For large $j$ put
\[
 c_j=\frac{\int\varphi\,d\sigma_j}{\int\eta\,d\sigma_j},
 \qquad \varphi_j=\varphi-c_j\eta.
\]
Then $c_j\to0$. In \eqref{eq:pairing-definition}, first restrict
both variables to a fixed large ball with null boundary for the limit
measure, and replace $K_n(x-y)$ by the continuous capped kernel
$(x-y)/\max\{h,|x-y|\}^{n+1}$. Weak convergence of product measures
gives convergence of the capped pairing on this bounded product.
The change on $|x-y|\le h$ is uniformly $O(h)$ by
\eqref{eq:near-growth-integral} and upper growth; the
exterior is uniformly $O(T^{-1})$ by Lemma~\ref{lem:pairings}.
The same estimates apply to the uniformly bounded corrected tests.
Letting $h\downarrow0$ and $T\to\infty$ proves the pairing assertion.
\end{proof}

For measures defined on larger and larger balls, the same argument is
local: each fixed support ball and test eventually lies in the region
where the hypotheses hold. Tangent measures below are limits of the
rescalings $F\mapsto r_j^{-n}\sigma(x+r_jF)$ with $r_j\downarrow0$.
They are nonzero by lower regularity. Upper and lower growth, bounded
smooth annuli, and the absence of a flat subball on fixed compact
regions pass to such limits with slightly relaxed constants. The last
assertion follows from the two support statements and finite nets;
one first asks for strict flatness in a slightly larger ball.

\subsection{Flatness and regular dyadic cells}

\begin{definition}\label{def:flatness}
For $r>0$ and a measure $\sigma$ with support $F$, let
\begin{align*}
 e_\sigma(a,r)
 &=\inf_L\left(r^{-n-2}\int_{B(a,r)}\dist(x,L)^2\,d\sigma(x)\right)^{1/2},\\
 b_\sigma(a,r)
 &=\inf_L\frac1r\left(
   \sup_{x\in F\cap B(a,r)}\dist(x,L)
   +\sup_{y\in L\cap B(a,r)}\dist(y,F)\right).
\end{align*}
Both infima run over affine $n$-planes in $\R^d$, and an empty
supremum is zero. We call $e$ the \emph{excess} and $b$ the
\emph{bilateral flatness number}.
\end{definition}

The second term in $b$ rules out holes in an approximating plane.
The excess measures only the support's squared distance from a plane.
We use approximate minimizers for either infimum; no attainment is
needed. Under a rescaling by $r^{-n}\sigma(a+r\,\cdot)$, both
quantities at $(a,rt)$ become the corresponding quantities at $(0,t)$.

We use the classical dyadic decomposition of an AD-regular set;
see \cite[Section~4, p.~12]{AzzamMourgoglouVilla2023} for its scale and
finite-diameter conventions, and \cite{DavidSemmes1993} for the
construction. We record exactly the properties needed.

\begin{proposition}[Regular dyadic cells]\label{prop:dyadic-lattice}
If $D_E>0$, there are a common $\mu$-null set $N\subset E$ and
families $\mathcal D_k$ of Borel cells in $E\setminus N$ such that:
\begin{enumerate}[label=(\roman*)]
\item each $\mathcal D_k$ partitions $E\setminus N$, and cells of
different generations are disjoint or nested;
\item writing $\ell(Q)=2^{-k}$ for $Q\in\mathcal D_k$, there are
centers $z_Q\in Q$ and constants $0<c_0\le1/2$, $C_0\ge2$ for which
\[
 (E\setminus N)\cap B(z_Q,c_0\ell(Q))\subset Q
 \subset E\cap B(z_Q,C_0\ell(Q)),\qquad
 c_1\ell(Q)^n\le\mu(Q)\le C_1\ell(Q)^n;
\]
\item every non-top cell has one parent, and every cell has finitely
many children. When
$D_E=\infty$, all integer generations are available. When
$D_E<\infty$, there is a largest scale $\ell_0$ with
$D_E\le\ell_0<2D_E$, and finitely many cells at that scale.
\end{enumerate}
All constants depend only on $d,n,C_{\rm AD}$.
\end{proposition}

In the bounded case, choose the integer $k_0$ with
$D_E\le2^{-k_0}<2D_E$. Discard original generations coarser than this
one. If the original largest scale is finer, insert the finitely many
intervening generations, each consisting of the single cell
$E\setminus N$. All added scales are comparable to $D_E$, and
$\mu(E)\asymp D_E^n$ by regularity and Lemma~\ref{lem:global-growth}.
Thus their core, outer-ball and mass bounds hold after adjusting the
constants. The original top generation is finite, so the finite-children
property also holds at the join. The scale and core constants can be
enlarged or decreased by fixed factors.
Removing the common null boundary set gives a simultaneous partition
at every generation. If necessary, move each center a sufficiently
small distance into the conull part of its core and decrease $c_0$;
this puts every center in its own cell. The mass bounds follow from
the cores, the outer balls, and Lemma~\ref{lem:global-growth}.
Small-boundary estimates are not needed below.

Write $\mathcal D$ for the generation-indexed union of the
$\mathcal D_k$. We also call its cells \emph{dyadic cubes}; no Euclidean
cubical shape is assumed. Cells at different generations remain distinct objects
even if their underlying sets coincide. A proper descendant means a
descendant at a strictly finer generation; containment notation for
cells in sums or chains includes this generation relation. A subfamily
$\mathcal F\subset\mathcal D$ is \emph{Carleson with constant $C$}
if
\begin{equation}\label{eq:carleson-definition}
 \sum_{Q\in\mathcal F,\ Q\subset P}\mu(Q)\le C\mu(P)
 \qquad(P\in\mathcal D).
\end{equation}
All sums are sums of nonnegative terms; proving the inequality for
every finite subfamily is sufficient. Each generation of descendants
partitions its parent up to the common null set. Consequently any
fixed number of generations contributes at most that number times
$\mu(P)$ to \eqref{eq:carleson-definition}.

For $J>1$ and $Q\in\mathcal D$, define
\begin{equation}\label{eq:oscillation-definition}
 W_J(Q)=\sup\left\{|\mathcal R_\mu(\varphi)|:
 \begin{array}{l}
 \varphi\in\Lip_c(\R^d),\ \supp\varphi\subset B(z_Q,J\ell(Q)),\\
 \Lip\varphi\le\ell(Q)^{-1},\quad\int\varphi\,d\mu=0
 \end{array}\right\}.
\end{equation}
Here $\Lip_c$ denotes compactly supported Lipschitz functions and
$\Lip\varphi$ their Lipschitz seminorm. The quantity $W_J(Q)$ has
the same scaling as $\mu(Q)$.

\subsection{Proof map}

Our geometric target is the Carleson estimate for
\[
 \{Q\in\mathcal D:b_\mu(z_Q,H\ell(Q))>\varepsilon\},
\]
for every $\varepsilon>0$ and a fixed sufficiently large $H$.
The David--Semmes bilateral weak geometric lemma will then supply
the Lipschitz images. The proof has three analytic stages and
one counting stage.

First, the operator bound makes the family
$W_J(Q)>v\ell(Q)^n$ Carleson for each fixed $J,v$.
It also supplies, outside a second Carleson family, a bilaterally flat
descendant within a fixed number of generations. These are
Propositions~\ref{prop:oscillation-packing} and~\ref{prop:flat-seeds}.

Second, a sequence whose excess and transform oscillation vanish has
a flat limiting measure. The plane comparison and reflectionless
arguments in Proposition~\ref{prop:flat-limit} identify that measure
as a positive constant times plane measure. This is the base measure
for the height argument; it does not yet give relative excess decay.

The affine-rigidity statement, Lemma~\ref{lem:fractional-rigidity}, is
independent of these measure limits. We use it twice: first, its bounded
case makes the density of a planar reflectionless measure constant in
Lemma~\ref{lem:planar-rigidity}; later, its full weighted form identifies
the normalized heights. Its proof in
Section~\ref{sec:fractional-rigidity} uses neither planar-measure rigidity
nor height compactness, so these two applications introduce no circular
dependence.

Third, Lemmas~\ref{lem:height-energy} and~\ref{lem:moving-compactness}
give compactness of heights divided by their initial excess scale.
Proposition~\ref{prop:height-equation} and
Lemma~\ref{lem:fractional-rigidity} identify the common limit as
affine. Theorem~\ref{thm:excess-propagation} converts that conclusion
into a uniform implication: small excess over a finite range of
larger scales, together with small $W_J$, controls excess and
bilateral flatness at the next smaller scale.

\begin{figure}[ht]
\centering
\begin{tikzpicture}[>=Stealth,
 box/.style={draw,rounded corners=2pt,align=center,text width=3.1cm,
 minimum height=1.2cm,font=\small},node distance=1.2cm]
 \node[box] (a) {Small excess\\and oscillation};
 \node[box,right=of a] (b) {Flat limiting\\measure};
 \node[box,right=of b] (c) {Heights divided by\\the excess scale};
 \node[box,below=1cm of c] (d) {Convergent moments\\and fractional equation};
 \node[box,left=of d] (e) {Affine height\\limit};
 \node[box,left=of e] (f) {Relative excess\\improvement};
 \draw[->] (a)--(b); \draw[->] (b)--(c); \draw[->] (c)--(d);
 \draw[->] (d)--(e); \draw[->] (e)--(f);
\end{tikzpicture}
\caption{The second limiting argument retains the size of the incoming
error. Its affine conclusion supplies the improvement needed to
propagate flatness. The independent weighted-rigidity lemma also
identifies the density in the flat limiting measure.}\label{fig:two-limits}
\end{figure}

Finally, start at a flat descendant and follow a chain of cells.
A later cell with large bilateral flatness must have an intervening
oscillation-bad cell. Lemma~\ref{lem:seed-charging} counts these
interventions for an arbitrary finite bad family. This proves
Proposition~\ref{prop:beta-packing}, and the geometric endpoint then
completes Theorem~\ref{thm:main}.
